% ===============================================================
% Document Class ================================================
\documentclass{article}

% ===============================================================
% Graphic Packages ==============================================
\usepackage{graphicx}       % Permite a inclusão de imagens no documento.
\usepackage{tocloft} % Pacote para customização de listas de conteúdos, figuras e tabelas
\usepackage{tikz}           % Ferramenta poderosa para criar gráficos programaticamente dentro do LaTeX.
\usetikzlibrary{calc}       % Extensão da biblioteca TikZ que permite cálculos mais complexos de coordenadas.
\usepackage[portuguese]{babel}
\usepackage{tocloft} % Pacote para customização de listas
\renewcommand{\listfigurename}{Lista de Figuras} % Altera o título para português
\AtBeginDocument{\renewcommand{\contentsname}{Sumário}} % Altera nome de conteudo para sumário
% ===============================================================
% Mathematical Tools ============================================
\usepackage{amsmath}        % Melhora a aparência e a flexibilidade de comandos matemáticos.
\usepackage{siunitx}        % Facilita o uso de unidades do Sistema Internacional e ajuda a formatar números complexos.

% ===============================================================
% Font and Text Appearance ======================================
\usepackage{mathptmx}       % Altera a fonte padrão do documento para Times New Roman.

% ===============================================================
% Table of Contents Customization ===============================
\usepackage{tocloft}  % Oferece controle total sobre a aparência das listas de conteúdos, figuras, tabelas, etc.

% ===============================================================
% Figure Positioning ============================================
\usepackage{float}     % Melhora a interface para definir o posicionamento de objetos flutuantes como figuras e tabelas.

% ===============================================================
% Paragraph Spacing and Indentation =============================
\usepackage{setspace}  % Permite o ajuste fino do espaçamento entre linhas.
\usepackage{indentfirst} % Adiciona indentação ao primeiro parágrafo de cada seção.

% ===============================================================
% Page Layout ===================================================
\usepackage[a4paper, top=3cm, bottom=2cm, left=3cm, right=2cm]{geometry}  % Define as margens de todo o documento.
\setlength{\parindent}{4em}  % Define o tamanho da indentação para todos os parágrafos.
\setlength{\emergencystretch}{3em}

% ===============================================================
% Section Heading Customization =================================
\makeatletter
\renewcommand\paragraph{\@startsection{paragraph}{4}{\z@}%
    {2ex plus 1ex minus .2ex}%
    {1em}%
    {\normalfont\normalsize\bfseries}}
\makeatother

% ===============================================================
% Urls ==========================================================
\usepackage{hyperref} % Para criar links clicáveis
\hypersetup{
    colorlinks=true,
    linkcolor=blue,
    filecolor=magenta,
    urlcolor=blue,
    citecolor=blue,
    pdfborder={0 0 0},  % Remove o quadrado ao redor dos links
    breaklinks=true  % Permite que os links sejam quebrados entre as linhas
}
\usepackage{xurl} % Melhora a quebra de links em qualquer lugar

% ===============================================================
% Bloco de código ===============================================
\usepackage{listings}
\usepackage{xcolor}
\usepackage{courier}
\lstset{
  backgroundcolor=\color[RGB]{249,246,239},
  basicstyle=\ttfamily\footnotesize,
  breaklines=true,
  frame=single,
  numbers=left,
  numberstyle=\tiny\color{gray},
  keywordstyle=\color[RGB]{40,40,255},
  commentstyle=\color[RGB]{0,125,0},
  stringstyle=\color[RGB]{255,0,0},
  showstringspaces=false,
  rulecolor=\color{black},
  captionpos=b,
  abovecaptionskip=5pt,
  belowcaptionskip=5pt,
  xleftmargin=0.15\textwidth,
  xrightmargin=0.15\textwidth,
  morecomment=[l]{//}
}
\renewcommand{\lstlistlistingname}{Lista de Códigos} % Comando para renomear o título da Lista de Códigos

% ===============================================================
% Begin Document ================================================
\begin{document}

% ===============================================================
% Capa ==========================================================
\begin{titlepage}
    \centering
    % Desenhar a margem
    \begin{tikzpicture}[remember picture, overlay]
        \draw[line width = 4pt] ($(current page.north west) + (20mm, -20mm)$) rectangle ($(current page.south east) + (-10mm,10mm)$);
        \draw[line width = 1pt] ($(current page.north west) + (21.5mm,-21.5mm)$) rectangle ($(current page.south east) + (-11.5mm,11.5mm)$);
    \end{tikzpicture}

    % Linha de logos
    \includegraphics[height=0.151\textwidth]{header/contra-capa/assets/uerj.png}\hfill
    \includegraphics[height=0.15\textwidth]{header/contra-capa/assets/iprj.jpeg}\hfill

    \vspace{2cm} % Espaço vertical após os logos

    % Informações do curso
    {\Large\bfseries Instituto Politécnico do Estado do Rio de Janeiro \par}
    \vspace{0.5cm}
    {\large Curso de Engenharia da Computação \par}

    \vspace{3cm} % Espaço vertical antes do título da atividade

    % Nome do estudante
    {\large\bfseries Guilherme Cagide Fialho \par}

    \vspace{3cm}

    % Título da atividade
    {\large\bfseries Trabalho 2 - Microsoft\par}

    \vfill % Empurra o restante para o fundo da página

    % Local e ano
    {\large\bfseries Nova Friburgo \par}
    \vspace{0.3cm}
    {\large\bfseries 2024 \par}
\end{titlepage}
\newpage % Começa uma nova página após a capa


% ===============================================================
% Contra Capa ===================================================
\begin{titlepage}
    \centering
    % Desenhar a margem
    \begin{tikzpicture}[remember picture, overlay]
        \draw[line width = 4pt] ($(current page.north west) + (20mm, -20mm)$) rectangle ($(current page.south east) + (-10mm,10mm)$);
        \draw[line width = 1pt] ($(current page.north west) + (21.5mm,-21.5mm)$) rectangle ($(current page.south east) + (-11.5mm,11.5mm)$);
    \end{tikzpicture}

    % Linha de logos
    \includegraphics[height=0.151\textwidth]{header/contra-capa/assets/uerj.png}\hfill
    \includegraphics[height=0.15\textwidth]{header/contra-capa/assets/iprj.jpeg}\hfill

    \vspace{2cm} % Espaço vertical após os logos

    % Informações do curso
    {\Large\bfseries Instituto Politécnico do Estado do Rio de Janeiro \par}
    \vspace{0.5cm}
    {\large Graduação em Engenharia da Computação \par}

    \vspace{3cm} % Espaço vertical antes do título da atividade

    % Nome do estudante
    {\large Guilherme Cagide Fialho \par}

    \vspace{1.5cm}

    % Título da atividade
    {\large\bfseries Trabalho 5 - Amazon \par}

    \vspace{1cm} % Espaço vertical antes do título da atividade

    % Informações do projeto
    \begin{flushright}
        \begin{minipage}{0.5\textwidth}
            \large
            \raggedleft % Garante que o texto dentro da minipage seja alinhado à direita
            Trabalho da Disciplina Estudo de Casos Empresariais
        \end{minipage}
    \end{flushright}

    \vspace{1.5cm}

    % Informações do orientador
    \begin{flushleft}
        \begin{minipage}{0.5\textwidth}
            \large
            \raggedright
            Orientador: Prof. Leandro de Amorim Ratamero
        \end{minipage}
    \end{flushleft}


    \vfill % Empurra o restante para o fundo da página

    % Local e ano
    {\large Nova Friburgo \par}
    \vspace{0.3cm}
    {\large 2024 \par}
\end{titlepage}
\newpage % Começa uma nova página após a capa


% ===============================================================
% Resumo ========================================================
% \begin{titlepage}
    \thispagestyle{empty} % Remove números de página
    \setstretch{1.5} % Espaçamento entre linhas, certifique-se de que o pacote setspace está incluído em document.tex

    \begin{center}
        \textbf{\Large RESUMO}
    \end{center}


    \vspace{1cm} % Espaço vertical

    \noindent CAGIDE FIALHO, G. Relatório do projeto de Modelagem
    e controle de sistemas. 2024. 54 f. Trabalho de Conclusão de Disciplina Modelagem e Controle de Sistemas (Graduação em
    Engenharia da computação) – Graduação em Engenharia da Computação, Universidade
    do Estado do Rio de Janeiro, Nova Friburgo, 2024.

    \vspace{0.4cm} % Espaço vertical

    Este trabalho explora a modelagem e o controle de um sistema dinâmico do tipo massa-mola-amortecedor, utilizando a plataforma Scilab para desenvolvimento e simulação. O foco do estudo está na implementação de modelos matemáticos para descrever a dinâmica do sistema e na análise de sua resposta sob diversas condições iniciais, sem a aplicação de forças externas. Utilizando a ferramenta Xcos, um componente gráfico do Scilab, realizamos simulações que permitiram uma análise visual e quantitativa das respostas transientes do sistema. O estudo destaca a influência dos parâmetros físicos, como a massa, o coeficiente de amortecimento, e a constante da mola, nas características de resposta do sistema. Além disso, técnicas de controle foram empregadas para ajustar a resposta do sistema, demonstrando como o amortecimento pode contribuir para a estabilização após perturbações e enfatizando a relevância de uma parametrização cuidadosa para alcançar um comportamento eficaz do sistema. Este projeto contribui para a compreensão das teorias de controle aplicáveis em sistemas mecânicos e outros contextos de sistemas dinâmicos na engenharia.
    \vspace{0.4cm} % Espaço vertical

    \textbf{Palavras-chave}: Modelagem e Controle, Sistema Massa-Mola-Amortecedor, Simulação, Scilab, Xcos.
\end{titlepage}


% ===============================================================
% Abstract ======================================================
% \begin{titlepage}
    \thispagestyle{empty} % Remove page numbers
    \setstretch{1.5} % Line spacing, make sure the setspace package is included in document.tex

    \begin{center}
        \textbf{\Large ABSTRACT}
    \end{center}

    \vspace{1cm} % Vertical space

    \noindent CAGIDE FIALHO, G. Report on the Modeling and Control Systems Project. 2024. 54 p. Final Project for the Modeling and Control Systems Course (Bachelor’s degree in Computer Engineering) - Bachelor’s Degree in Computer Engineering, State University of Rio de Janeiro, Nova Friburgo, 2024.

    \vspace{0.4cm} % Vertical space

    This work explores the modeling and control of a dynamic system of the mass-spring-damper type, using the Scilab platform for development and simulation. The focus of the study is on implementing mathematical models to describe the system dynamics and analyzing its response under various initial conditions without the application of external forces. Using Xcos, a graphical component of Scilab, we conducted simulations that allowed for visual and quantitative analysis of the system's transient responses. The study highlights the influence of physical parameters, such as mass, damping coefficient, and spring constant, on the system's response characteristics. Furthermore, control techniques were employed to adjust the system's response, demonstrating how damping can contribute to stabilization after disturbances and emphasizing the importance of careful parameterization to achieve effective system behavior. This project contributes to the understanding of control theories applicable to mechanical systems and other dynamic system contexts in engineering.
    \vspace{0.4cm} % Vertical space

    \textbf{Keywords}: Modeling and Control, Mass-Spring-Damper System, Simulation, Scilab, Xcos.
\end{titlepage}


% ===============================================================
% Lista de figuras ==============================================
% \newpage
\setstretch{1.5} % Espaçamento entre linhas, certifique-se de que o pacote setspace está incluído em document.tex

% % Estilo de citação (plain, abbrv, ieeetr, etc.)
% \bibliographystyle{IEEEtran}  % Estilo que inclui URLs
% % \bibliographystyle{plain}

% \bibliography{
%     perguntas/1-foundation/latex/references,
%     perguntas/2-graduation/latex/references,
%     perguntas/3-main-productions/latex/references,
%     perguntas/4-location/latex/references,
%     perguntas/5-competitors/latex/references,
%     perguntas/6-protection/latex/references
% }

% Estilo de citação (plain, abbrv, ieeetr, etc.)
\bibliographystyle{IEEEtran}  % Estilo que inclui URLs
% \bibliographystyle{plain}

\bibliography{perguntas/1-foundation/latex/references,perguntas/2-graduation/latex/references,perguntas/3-main-products/latex/references.bib,perguntas/4-location/latex/references,perguntas/5-competitors/latex/references,perguntas/6-protection/latex/references, perguntas/7-profit/latex/references, perguntas/8-train/latex/references, perguntas/9-social-programs/latex/references, perguntas/10-green-purposes/latex/references}


% ===============================================================
% Lista de códigos ==============================================
% \newpage
\setstretch{1.5} % Espaçamento entre linhas, certifique-se de que o pacote setspace está incluído em document.tex

% % Estilo de citação (plain, abbrv, ieeetr, etc.)
% \bibliographystyle{IEEEtran}  % Estilo que inclui URLs
% % \bibliographystyle{plain}

% \bibliography{
%     perguntas/1-foundation/latex/references,
%     perguntas/2-graduation/latex/references,
%     perguntas/3-main-productions/latex/references,
%     perguntas/4-location/latex/references,
%     perguntas/5-competitors/latex/references,
%     perguntas/6-protection/latex/references
% }

% Estilo de citação (plain, abbrv, ieeetr, etc.)
\bibliographystyle{IEEEtran}  % Estilo que inclui URLs
% \bibliographystyle{plain}

\bibliography{perguntas/1-foundation/latex/references,perguntas/2-graduation/latex/references,perguntas/3-main-products/latex/references.bib,perguntas/4-location/latex/references,perguntas/5-competitors/latex/references,perguntas/6-protection/latex/references, perguntas/7-profit/latex/references, perguntas/8-train/latex/references, perguntas/9-social-programs/latex/references, perguntas/10-green-purposes/latex/references}


% ===============================================================
% Lista de siglas e abreviaturas ================================
% \newpage
\setstretch{1.5} % Espaçamento entre linhas, certifique-se de que o pacote setspace está incluído em document.tex

% % Estilo de citação (plain, abbrv, ieeetr, etc.)
% \bibliographystyle{IEEEtran}  % Estilo que inclui URLs
% % \bibliographystyle{plain}

% \bibliography{
%     perguntas/1-foundation/latex/references,
%     perguntas/2-graduation/latex/references,
%     perguntas/3-main-productions/latex/references,
%     perguntas/4-location/latex/references,
%     perguntas/5-competitors/latex/references,
%     perguntas/6-protection/latex/references
% }

% Estilo de citação (plain, abbrv, ieeetr, etc.)
\bibliographystyle{IEEEtran}  % Estilo que inclui URLs
% \bibliographystyle{plain}

\bibliography{perguntas/1-foundation/latex/references,perguntas/2-graduation/latex/references,perguntas/3-main-products/latex/references.bib,perguntas/4-location/latex/references,perguntas/5-competitors/latex/references,perguntas/6-protection/latex/references, perguntas/7-profit/latex/references, perguntas/8-train/latex/references, perguntas/9-social-programs/latex/references, perguntas/10-green-purposes/latex/references}


% ===============================================================
% Materiais utilizados ==========================================
% \newpage
\setstretch{1.5} % Espaçamento entre linhas, certifique-se de que o pacote setspace está incluído em document.tex

% % Estilo de citação (plain, abbrv, ieeetr, etc.)
% \bibliographystyle{IEEEtran}  % Estilo que inclui URLs
% % \bibliographystyle{plain}

% \bibliography{
%     perguntas/1-foundation/latex/references,
%     perguntas/2-graduation/latex/references,
%     perguntas/3-main-productions/latex/references,
%     perguntas/4-location/latex/references,
%     perguntas/5-competitors/latex/references,
%     perguntas/6-protection/latex/references
% }

% Estilo de citação (plain, abbrv, ieeetr, etc.)
\bibliographystyle{IEEEtran}  % Estilo que inclui URLs
% \bibliographystyle{plain}

\bibliography{perguntas/1-foundation/latex/references,perguntas/2-graduation/latex/references,perguntas/3-main-products/latex/references.bib,perguntas/4-location/latex/references,perguntas/5-competitors/latex/references,perguntas/6-protection/latex/references, perguntas/7-profit/latex/references, perguntas/8-train/latex/references, perguntas/9-social-programs/latex/references, perguntas/10-green-purposes/latex/references}


% ===============================================================
% Sumário =======================================================
\newpage
\setstretch{1.5} % Espaçamento entre linhas, certifique-se de que o pacote setspace está incluído em document.tex

% % Estilo de citação (plain, abbrv, ieeetr, etc.)
% \bibliographystyle{IEEEtran}  % Estilo que inclui URLs
% % \bibliographystyle{plain}

% \bibliography{
%     perguntas/1-foundation/latex/references,
%     perguntas/2-graduation/latex/references,
%     perguntas/3-main-productions/latex/references,
%     perguntas/4-location/latex/references,
%     perguntas/5-competitors/latex/references,
%     perguntas/6-protection/latex/references
% }

% Estilo de citação (plain, abbrv, ieeetr, etc.)
\bibliographystyle{IEEEtran}  % Estilo que inclui URLs
% \bibliographystyle{plain}

\bibliography{perguntas/1-foundation/latex/references,perguntas/2-graduation/latex/references,perguntas/3-main-products/latex/references.bib,perguntas/4-location/latex/references,perguntas/5-competitors/latex/references,perguntas/6-protection/latex/references, perguntas/7-profit/latex/references, perguntas/8-train/latex/references, perguntas/9-social-programs/latex/references, perguntas/10-green-purposes/latex/references}


% ===============================================================
% Introdução ====================================================
% \newpage
\setstretch{1.5} % Espaçamento entre linhas, certifique-se de que o pacote setspace está incluído em document.tex

% % Estilo de citação (plain, abbrv, ieeetr, etc.)
% \bibliographystyle{IEEEtran}  % Estilo que inclui URLs
% % \bibliographystyle{plain}

% \bibliography{
%     perguntas/1-foundation/latex/references,
%     perguntas/2-graduation/latex/references,
%     perguntas/3-main-productions/latex/references,
%     perguntas/4-location/latex/references,
%     perguntas/5-competitors/latex/references,
%     perguntas/6-protection/latex/references
% }

% Estilo de citação (plain, abbrv, ieeetr, etc.)
\bibliographystyle{IEEEtran}  % Estilo que inclui URLs
% \bibliographystyle{plain}

\bibliography{perguntas/1-foundation/latex/references,perguntas/2-graduation/latex/references,perguntas/3-main-products/latex/references.bib,perguntas/4-location/latex/references,perguntas/5-competitors/latex/references,perguntas/6-protection/latex/references, perguntas/7-profit/latex/references, perguntas/8-train/latex/references, perguntas/9-social-programs/latex/references, perguntas/10-green-purposes/latex/references}


% ===============================================================
% Atividades ====================================================
\section{MEI, Microempreendedor Individual}

O Microempreendedor Individual (MEI) constitui uma categoria empresarial no Brasil projetada para regularizar a atuação de pequenos empreendedores que, frequentemente, operam de maneira autônoma e informal. Essa categoria facilita a formalização desses empreendedores como empresas através de um processo simplificado e econômico. O objetivo principal da implementação do MEI é promover a inclusão desses trabalhadores autônomos no sistema econômico formal, proporcionando uma estrutura que permite a transição do informal para o formal com mínimas barreiras burocráticas e financeiras.

\subsection{Características Principais}

Algumas das principais características do MEI incluem:
\begin{itemize}
    \item Faturamento máximo permitido de até R\$ 81.000,00 por ano.
    \item Não é permitida a participação em outra empresa como sócio ou titular.
    \item O MEI pode ter no máximo um empregado contratado que receba o salário mínimo ou o piso da categoria.
    \item É obrigatório contribuir mensalmente para o INSS, garantindo assim benefícios como aposentadoria e auxílio-doença.
\end{itemize}

\subsection{Benefícios}

Entre os benefícios de ser um MEI, estão:
\begin{itemize}
    \item Emissão de nota fiscal eletrônica.
    \item Acesso a serviços bancários, incluindo crédito.
    \item Baixo custo de tributação, com valores fixos mensais que cobrem tributos como INSS e ISS ou ICMS, dependendo do tipo de serviço ou produto fornecido.
    \item Possibilidade de participar de licitações públicas.
    \item O pagamento do imposto mensal, que em 2024 está em cerca de R\$ 71,60, é considerado baixo em comparação com outros regimes tributários e já conta como tempo de contribuição para a aposentadoria.
\end{itemize}

\subsection{Obrigações}

As obrigações do MEI incluem:
\begin{itemize}
    \item Realização de uma declaração anual simplificada.
    \item Manutenção de um registro das vendas e compras realizadas.
    \item Emissão de nota fiscal quando o cliente for outra empresa.
\end{itemize} % Completa e Validada
\section{EI, Empresa Individual}

A Empresa Individual (EI) é uma forma de organização empresarial que permite ao empreendedor atuar como único proprietário, sem a necessidade de formar sociedade. Diferentemente do MEI, a Empresa Individual permite um faturamento maior e não impõe limite de contratação de empregados, mas expõe o empresário a uma maior responsabilidade financeira, uma vez que não há distinção entre o patrimônio pessoal e o empresarial.

\subsection{Características Principais}

As principais características da Empresa Individual incluem:
\begin{itemize}
    \item Não há limite de faturamento anual imposto pela legislação específica do EI.
    \item O empresário responde ilimitadamente pelas dívidas da empresa, ou seja, seu patrimônio pessoal pode ser utilizado para quitações de débitos.
    \item Flexibilidade na gestão e operação do negócio, sem a necessidade de consultas ou aprovações de sócios.
\end{itemize}

\subsection{Benefícios}

Os benefícios de optar por uma Empresa Individual incluem:
\begin{itemize}
    \item Autonomia total na gestão empresarial, com decisões rápidas e sem necessidade de deliberações societárias.
    \item Processo de abertura e encerramento do negócio geralmente mais simples e menos burocrático em comparação com outras formas societárias.
    \item Menos exigências legais e contábeis do que uma sociedade limitada ou uma sociedade anônima.
\end{itemize}

\subsection{Obrigações}

As obrigações de uma Empresa Individual são significativas, especialmente no que tange à responsabilidade por dívidas:
\begin{itemize}
    \item Manutenção de registros financeiros e contábeis atualizados.
    \item Cumprimento de todas as obrigações fiscais e tributárias aplicáveis, dependendo do regime tributário adotado (Simples Nacional, Lucro Presumido ou Lucro Real).
    \item Emissão regular de notas fiscais para todas as transações comerciais.
\end{itemize}
 % Completa e Validadada
\section{SLU, Sociedade Limitada Unipessoal}

A Sociedade Limitada Unipessoal (SLU) é uma forma jurídica que permite a constituição de uma empresa limitada por apenas um sócio, sem exigência de capital social mínimo. Esta modalidade foi introduzida para facilitar a formalização de empreendedores individuais, oferecendo proteção patrimonial sem necessidade de associação com outros sócios.

\subsection{Características Principais}

As características da SLU incluem:
\begin{itemize}
    \item Proteção do patrimônio pessoal do proprietário, diferentemente do EI, onde o patrimônio pessoal e empresarial são indistintos.
    \item Não requer um capital social mínimo para sua constituição, ao contrário da EIRELI, que exige um capital correspondente a 100 salários mínimos.
    \item Flexibilidade para operar em qualquer setor econômico, incluindo profissões regulamentadas como médicos, engenheiros e advogados.
\end{itemize}

\subsection{Vantagens}

Entre as vantagens da SLU, estão:
\begin{itemize}
    \item Facilidade de gestão com autonomia total nas decisões empresariais, sem necessidade de deliberação com outros sócios.
    \item Simplificação do processo de abertura e gestão empresarial, com menos burocracia comparada a outras formas jurídicas.
    \item Permissão para estabelecer mais de uma empresa nessa modalidade, aumentando as possibilidades de expansão de negócios para o empreendedor.
\end{itemize}

\subsection{Processo de Abertura}

O processo de abertura de uma SLU é relativamente simples e inclui:
\begin{itemize}
    \item Registro do contrato social na Junta Comercial do estado.
    \item Inscrição no CNPJ através do site da Receita Federal.
    \item Solicitação de alvará de funcionamento junto à prefeitura local.
\end{itemize}

\subsection{Considerações Finais}

A SLU representa uma opção vantajosa para empreendedores que desejam a segurança de uma sociedade limitada com a simplicidade de gestão de um negócio unipessoal. É uma escolha popular entre os pequenos empresários que buscam uma entrada menos complicada no mundo empresarial.
 % Completa e validada

\section{LTDA, Sociedade Limitada Empresarial}

A Sociedade Limitada (LTDA) é uma das formas societárias mais comuns no Brasil, caracterizando-se pela limitação da responsabilidade dos sócios ao montante de suas cotas no capital social. Este modelo é preferido por pequenas e médias empresas devido à sua flexibilidade operacional e proteção patrimonial.

\subsection{Características Principais}
\begin{itemize}
    \item \textbf{Responsabilidade Limitada:} Os sócios têm sua responsabilidade limitada ao valor de suas cotas, protegendo os bens pessoais contra dívidas da empresa.
    \item \textbf{Capital Social:} Dividido em cotas, cada sócio contribui com dinheiro, bens ou direitos, valorizados monetariamente.
    \item \textbf{Administração:} Pode ser feita por sócios ou terceiros designados no contrato social.
    \item \textbf{Contrato Social:} Documento fundamental que regula a participação de cada sócio, distribuição de lucros, perdas e as regras de transferência de cotas.
\end{itemize}

\subsection{Vantagens}
\begin{itemize}
    \item \textbf{Proteção Patrimonial:} Separação entre patrimônio pessoal dos sócios e da empresa.
    \item \textbf{Flexibilidade Tributária:} Possibilidade de escolha entre Simples Nacional, Lucro Presumido ou Lucro Real.
    \item \textbf{Simplicidade Gerencial:} Menos burocracia na gestão, ideal para negócios familiares e de pequeno porte.
\end{itemize}

\subsection{Desvantagens}
\begin{itemize}
    \item \textbf{Conflitos Internos:} Potenciais conflitos na gestão podem surgir sem clareza no contrato social.
    \item \textbf{Rigidez para Alterações:} Mudanças estruturais requerem formalidades e consenso entre os sócios.
\end{itemize}

\subsection{Processo de Abertura}
\noindent Para abrir uma LTDA, são necessários os seguintes passos:
\begin{enumerate}
    \item Definir os sócios e estabelecer o capital social.
    \item Elaborar e registrar o contrato social na Junta Comercial.
    \item Obter o CNPJ e as licenças necessárias para operação.
\end{enumerate}

\subsection{Conclusão}
A Sociedade Limitada oferece uma estrutura flexível e segura para empreendedores que buscam proteção patrimonial e facilidade de gestão. É uma escolha robusta para quem deseja começar um negócio com riscos limitados aos investimentos realizados.
 % Completa e validada
\section{SS, Sociedade Limitada Simples}
A Sociedade Simples é uma forma de associação entre profissionais que exercem predominantemente atividades de natureza intelectual, científica, literária ou artística, com o propósito de oferecer serviços pessoais especializados. Este tipo de sociedade é regido principalmente pelos artigos 982 e 983 do Código Civil Brasileiro de 2002.

\subsection{Características Principais}
\begin{itemize}
    \item \textbf{Natureza não empresarial}: As atividades realizadas têm um caráter mais pessoal, direto e especializado, sem a necessidade de uma estrutura empresarial.
    \item \textbf{Responsabilidade dos sócios}: Dependendo da configuração, os sócios podem ter responsabilidade ilimitada, onde as dívidas sociais podem se estender aos bens pessoais.
    \item \textbf{Administração}: Deve ser realizada por pessoas naturais, com a administração detalhada no contrato social.
    \item \textbf{Registro}: Realizado no Registro Civil de Pessoas Jurídicas, diferentemente das sociedades empresariais que são registradas na Junta Comercial.
\end{itemize}

\subsection{Vantagens e Desvantagens}
\subsubsection{Vantagens}
\begin{itemize}
    \item Flexibilidade para profissionais que desejam colaborar sem complexidades empresariais.
    \item Permite contribuição dos sócios em forma de serviços, além de capital.
\end{itemize}

\subsubsection{Desvantagens}
\begin{itemize}
    \item Maior risco financeiro para os sócios devido à possível responsabilidade ilimitada.
    \item Menos proteção em relação à separação dos bens pessoais dos bens da sociedade.
\end{itemize}

\subsection{Processo de Abertura}
\begin{enumerate}
    \item Elaboração de um contrato social com detalhes claros sobre gestão, responsabilidades e direitos dos sócios.
    \item Registro do contrato no Registro Civil de Pessoas Jurídicas para obtenção do CNPJ.
    \item Obtenção de licenças necessárias para operação conforme as regulamentações locais.
\end{enumerate}

\subsection{Conclusão}
A Sociedade Simples é uma escolha ideal para profissionais que buscam uma cooperação estruturada com flexibilidade e uma abordagem pessoal na prestação de serviços. Embora ofereça vantagens significativas em termos de flexibilidade operacional, é importante considerar as implicações legais e financeiras associadas à sua estrutura de responsabilidade.
 % Completa e validada

\section{SA, Sociedade Anônima}
A Sociedade Anônima (S.A.) é uma forma de organização empresarial onde o capital social é dividido em ações, e a responsabilidade dos acionistas é limitada ao preço das ações que possuem. Este modelo é adequado tanto para grandes operações de mercado quanto para empreendimentos menores que preferem não negociar publicamente suas ações.

\subsection{Características Principais}
\begin{itemize}
    \item \textbf{Capital Dividido em Ações}: As participações são divididas em ações, que podem ser de dois tipos principais: ordinárias (com direito a voto) e preferenciais (sem direito a voto, mas com prioridade em dividendos).
    \item \textbf{Responsabilidade Limitada}: Os acionistas têm sua responsabilidade limitada ao valor de suas ações, protegendo seus bens pessoais de responsabilidades corporativas.
\end{itemize}

\subsection{Estrutura Organizacional}
\subsubsection{Órgãos Administrativos}
\begin{itemize}
    \item \textbf{Assembleia Geral}: Órgão máximo, responsável pelas decisões cruciais, incluindo eleição dos membros do conselho e alterações estatutárias.
    \item \textbf{Conselho de Administração}: Elege e supervisiona a diretoria, define estratégias e políticas corporativas.
    \item \textbf{Diretoria}: Responsável pela gestão diária e operacional da empresa.
    \item \textbf{Conselho Fiscal}: Supervisiona as operações financeiras e contábeis, assegurando a conformidade com a lei e o estatuto.
\end{itemize}

\subsection{Tipos de Sociedades Anônimas}
\begin{itemize}
    \item \textbf{Capital Aberto}: Ações negociadas em bolsas de valores, sujeitas a regras de transparência e governança rigorosas.
    \item \textbf{Capital Fechado}: Ações não disponíveis ao público geral, com menor regulamentação.
\end{itemize}

\subsection{Processo de Abertura}
A formação de uma S.A. exige a elaboração de um estatuto social, registro na Junta Comercial, e, para as de capital aberto, registro na Comissão de Valores Mobiliários (CVM).

\subsection{Conclusão}
As Sociedades Anônimas representam uma estrutura robusta para empresas que buscam capital e crescimento, com uma forte ênfase na separação entre gestão corporativa e propriedade acionária.
 % Completa e validada
\begin{titlepage}
    \thispagestyle{empty} % Remove números de página
    \setstretch{1.5} % Espaçamento entre linhas, certifique-se de que o pacote setspace está incluído em document.tex

    \begin{center}
        \textbf{\Large RESUMO}
    \end{center}


    \vspace{1cm} % Espaço vertical

    \noindent CAGIDE FIALHO, G. Relatório do projeto de Modelagem
    e controle de sistemas. 2024. 54 f. Trabalho de Conclusão de Disciplina Modelagem e Controle de Sistemas (Graduação em
    Engenharia da computação) – Graduação em Engenharia da Computação, Universidade
    do Estado do Rio de Janeiro, Nova Friburgo, 2024.

    \vspace{0.4cm} % Espaço vertical

    Este trabalho explora a modelagem e o controle de um sistema dinâmico do tipo massa-mola-amortecedor, utilizando a plataforma Scilab para desenvolvimento e simulação. O foco do estudo está na implementação de modelos matemáticos para descrever a dinâmica do sistema e na análise de sua resposta sob diversas condições iniciais, sem a aplicação de forças externas. Utilizando a ferramenta Xcos, um componente gráfico do Scilab, realizamos simulações que permitiram uma análise visual e quantitativa das respostas transientes do sistema. O estudo destaca a influência dos parâmetros físicos, como a massa, o coeficiente de amortecimento, e a constante da mola, nas características de resposta do sistema. Além disso, técnicas de controle foram empregadas para ajustar a resposta do sistema, demonstrando como o amortecimento pode contribuir para a estabilização após perturbações e enfatizando a relevância de uma parametrização cuidadosa para alcançar um comportamento eficaz do sistema. Este projeto contribui para a compreensão das teorias de controle aplicáveis em sistemas mecânicos e outros contextos de sistemas dinâmicos na engenharia.
    \vspace{0.4cm} % Espaço vertical

    \textbf{Palavras-chave}: Modelagem e Controle, Sistema Massa-Mola-Amortecedor, Simulação, Scilab, Xcos.
\end{titlepage}
 % Completo e Validado

% ===============================================================
% Conclusão geral ===============================================
% \newpage
\setstretch{1.5} % Espaçamento entre linhas, certifique-se de que o pacote setspace está incluído em document.tex

% % Estilo de citação (plain, abbrv, ieeetr, etc.)
% \bibliographystyle{IEEEtran}  % Estilo que inclui URLs
% % \bibliographystyle{plain}

% \bibliography{
%     perguntas/1-foundation/latex/references,
%     perguntas/2-graduation/latex/references,
%     perguntas/3-main-productions/latex/references,
%     perguntas/4-location/latex/references,
%     perguntas/5-competitors/latex/references,
%     perguntas/6-protection/latex/references
% }

% Estilo de citação (plain, abbrv, ieeetr, etc.)
\bibliographystyle{IEEEtran}  % Estilo que inclui URLs
% \bibliographystyle{plain}

\bibliography{perguntas/1-foundation/latex/references,perguntas/2-graduation/latex/references,perguntas/3-main-products/latex/references.bib,perguntas/4-location/latex/references,perguntas/5-competitors/latex/references,perguntas/6-protection/latex/references, perguntas/7-profit/latex/references, perguntas/8-train/latex/references, perguntas/9-social-programs/latex/references, perguntas/10-green-purposes/latex/references}


% ===============================================================
% Referencias ===================================================
\newpage
\setstretch{1.5} % Espaçamento entre linhas, certifique-se de que o pacote setspace está incluído em document.tex

% % Estilo de citação (plain, abbrv, ieeetr, etc.)
% \bibliographystyle{IEEEtran}  % Estilo que inclui URLs
% % \bibliographystyle{plain}

% \bibliography{
%     perguntas/1-foundation/latex/references,
%     perguntas/2-graduation/latex/references,
%     perguntas/3-main-productions/latex/references,
%     perguntas/4-location/latex/references,
%     perguntas/5-competitors/latex/references,
%     perguntas/6-protection/latex/references
% }

% Estilo de citação (plain, abbrv, ieeetr, etc.)
\bibliographystyle{IEEEtran}  % Estilo que inclui URLs
% \bibliographystyle{plain}

\bibliography{perguntas/1-foundation/latex/references,perguntas/2-graduation/latex/references,perguntas/3-main-products/latex/references.bib,perguntas/4-location/latex/references,perguntas/5-competitors/latex/references,perguntas/6-protection/latex/references, perguntas/7-profit/latex/references, perguntas/8-train/latex/references, perguntas/9-social-programs/latex/references, perguntas/10-green-purposes/latex/references}


% ===============================================================
% Anexo ou Apendices ============================================
% \newpage
\setstretch{1.5} % Espaçamento entre linhas, certifique-se de que o pacote setspace está incluído em document.tex

% % Estilo de citação (plain, abbrv, ieeetr, etc.)
% \bibliographystyle{IEEEtran}  % Estilo que inclui URLs
% % \bibliographystyle{plain}

% \bibliography{
%     perguntas/1-foundation/latex/references,
%     perguntas/2-graduation/latex/references,
%     perguntas/3-main-productions/latex/references,
%     perguntas/4-location/latex/references,
%     perguntas/5-competitors/latex/references,
%     perguntas/6-protection/latex/references
% }

% Estilo de citação (plain, abbrv, ieeetr, etc.)
\bibliographystyle{IEEEtran}  % Estilo que inclui URLs
% \bibliographystyle{plain}

\bibliography{perguntas/1-foundation/latex/references,perguntas/2-graduation/latex/references,perguntas/3-main-products/latex/references.bib,perguntas/4-location/latex/references,perguntas/5-competitors/latex/references,perguntas/6-protection/latex/references, perguntas/7-profit/latex/references, perguntas/8-train/latex/references, perguntas/9-social-programs/latex/references, perguntas/10-green-purposes/latex/references}


% ===============================================================
% End Document ==================================================
\end{document}

